%%=====================================================================
%% VERSAO 2 (revisao pos-SBESC 2026, #33972 rejeitado em 24/09/2026).
%% VERSAO EM PORTUGUES. Template: IEEE two-column conference.
%% A versao submetida permanece intacta em ../artigo/.
%% Rastreabilidade revisor -> alteracao: RASTREABILIDADE-REVISOES.md
%%
%% Resultados novos (faixas, modelos isolados, pesos, CV aninhada,
%% proporcao, gerador corrigido, evasao): ../experimentos-revisao.py ->
%% ../artefatos-modelos/experimentos-revisao.json
%%=====================================================================
\documentclass[conference]{IEEEtran}
\IEEEoverridecommandlockouts

\usepackage[utf8]{inputenc}
\usepackage[T1]{fontenc}
\usepackage[brazil]{babel}
\usepackage{cite}
\usepackage{amsmath,amssymb,amsfonts}
\usepackage{graphicx}
\usepackage{booktabs}
\usepackage{array}
\usepackage{tikz}
\usetikzlibrary{positioning,arrows.meta}
\usepackage[hidelinks]{hyperref}
\usepackage{url}
\usepackage{balance} % equilibra as duas colunas na última página (evita sobra em branco)

\begin{document}

\title{Autenticação Contextual Autônoma para Aplicações\\ Web Combinando um Ensemble de IA e\\ Auditoria em Blockchain Permissionada}

\author{\IEEEauthorblockN{João Guedes de Brito}
\IEEEauthorblockA{\textit{Programa de Pós-Graduação em Engenharia de Sistemas e Produtos (PPGESP)}\\
\textit{Instituto Federal da Bahia (IFBA)}\\
Salvador, Bahia, Brasil \\
joaoguedesdebrito@gmail.com}
\and
\IEEEauthorblockN{Cleber Jorge Lira de Santana}
\IEEEauthorblockA{\textit{Programa de Pós-Graduação em Engenharia de Sistemas e Produtos (PPGESP)}\\
\textit{Instituto Federal da Bahia (IFBA)}\\
Salvador, Bahia, Brasil \\
cleberlira@ifba.edu.br}
}

\maketitle

\begin{abstract}
Mecanismos tradicionais de autenticação baseados em senhas estáticas ou
em autenticação multifator (MFA) fixa permanecem amplamente adotados,
mas ignoram o contexto comportamental e ambiental de cada tentativa de
acesso e são cada vez mais explorados por ataques de \emph{credential
stuffing}, phishing e tomada de conta. Este artigo apresenta um sistema
autônomo de autenticação contextual para aplicações web que combina um
ensemble de aprendizado de máquina com auditoria em blockchain
permissionada. Cada login é descrito por quatorze critérios contextuais
e comportamentais, avaliados por um ensemble de Isolation Forest, Random
Forest e uma rede neural multicamada; o escore de risco resultante
conduz uma decisão adaptativa em três faixas (permitir, reforçar com MFA
ou bloquear). Propõe-se ainda uma camada de auditoria em que cada
decisão, com o escore e os critérios que a justificaram, é registrada em
uma rede Hyperledger Fabric; discute-se quando essa escolha se justifica
frente a logs assinados e bancos append-only. Nesta fase, a camada de
auditoria é um projeto: o protótipo persiste as decisões em um banco
relacional e a integração com o Fabric ainda não foi realizada, de modo
que nenhuma propriedade de imutabilidade é reivindicada. Em avaliação
\emph{offline} sobre um conjunto de
dados sintéticos com cerca de 50.000 registros de login (95\% legítimos,
5\% suspeitos), o ensemble alcançou 98,1\% de acurácia, 80,2\% de
F1-score e AUC-ROC de 0,99 no conjunto de teste, com taxa de falsos
positivos de 0,8\% no limiar binário escolhido e latência de inferência
de aproximadamente 33\,ms por requisição; com as três faixas, 7,3\% dos
ataques passam sem desafio e 2,7\% dos acessos legítimos sofrem atrito.
O ensemble supera cada modelo isolado em cerca de 3 pontos de F1, mas um
atacante que imita apenas os critérios de horário e ritmo reduz o recall
de 77,6\% para 50,7\%, e com um proxy residencial, para 18,0\%. Os rótulos provêm do próprio
gerador sintético; os resultados caracterizam, portanto, a viabilidade
do pipeline, e não desempenho esperado em produção.
\end{abstract}

\begin{IEEEkeywords}
autenticação contextual, autenticação baseada em risco, aprendizado de
máquina, ensemble, detecção de anomalias, blockchain, trilha de
auditoria, segurança da informação.
\end{IEEEkeywords}

\section{Introdução}
A autenticação é a primeira linha de defesa de aplicações web.
Historicamente, apoiou-se em um fator de conhecimento --- a senha ---,
cujas fragilidades são amplamente documentadas: senhas são reutilizadas
entre serviços, vazadas em larga escala e exploradas por força bruta,
\emph{credential stuffing} e phishing~\cite{nist80063b,verizon2024}. A
autenticação multifator (MFA) mitiga parte dessa exposição ao combinar
fatores de conhecimento, posse e inerência, mas a MFA estática ainda
sucumbe a engenharia social, troca de chip (SIM swapping) e interceptação
de códigos em tempo real, além de introduzir atrito na experiência do
usuário~\cite{nist80063b}.

A escala do problema justifica ir além da credencial. Credential
stuffing, phishing e o uso de credenciais roubadas seguem entre os
principais vetores em relatórios de incidentes, e a fraude por tomada de
conta impõe custos financeiros e reputacionais diretos a organizações e
usuários finais~\cite{verizon2024}. A reutilização de senhas transforma
qualquer vazamento isolado em um passivo entre serviços, enquanto
ferramentas de ataque comoditizadas tornam o abuso de login em larga
escala barato e altamente automatizado. Sobretudo, defesas estáticas
--- políticas de robustez de senha e códigos de uso único --- não
observam \emph{como}, \emph{quando} ou \emph{de onde} ocorre um login e,
portanto, não distinguem um usuário legítimo de um adversário que já
detém credenciais válidas. Marcos regulatórios como a
LGPD~\cite{lgpd2018} acrescentam um requisito adicional: decisões
automatizadas de acesso sobre dados pessoais precisam ser comprovadamente
auditáveis. Em conjunto, essas pressões motivam decisões de autenticação
condicionadas ao contexto completo de cada tentativa --- e que deixem
uma trilha verificável do \emph{porquê} de cada decisão.

Para enfrentar essas limitações, as diretrizes de identidade digital do
NIST e o modelo de Zero Trust recomendam mecanismos \emph{adaptativos},
nos quais o nível de verificação exigido é proporcional ao risco
estimado de cada tentativa~\cite{nist80063b,rose2020}. Isso desloca a
autenticação de uma decisão binária e estática para uma avaliação
contínua e sensível ao contexto --- abordagem conhecida na literatura
como autenticação baseada em risco (\emph{Risk-Based Authentication},
RBA). Estudos empíricos de serviços em produção mostram, contudo, que o
RBA implantado é dominado por um único sinal (tipicamente o endereço IP)
e por heurísticas ajustadas manualmente~\cite{wiefling2019,wiefling2023}.
Em paralelo, quando as decisões de acesso são delegadas a modelos
autônomos, elas precisam ser \emph{auditáveis} por terceiros (por
exemplo, em fiscalizações de conformidade), requisito que registros de
decisão sob a guarda de um único operador não atendem plenamente.

Este artigo propõe um sistema autônomo de autenticação contextual que
acopla uma decisão adaptativa por aprendizado de máquina (ML) ao seu
registro para auditoria. As principais contribuições são:
\begin{itemize}
  \item uma arquitetura ponta a ponta que funde quatorze critérios
        contextuais e comportamentais em um único escore de risco
        calculado por um ensemble de Isolation Forest, Random Forest e
        rede neural, conduzindo uma decisão adaptativa em três faixas;
  \item o projeto de uma camada de auditoria em blockchain
        permissionada acoplada à decisão, com justificativa explícita
        de \emph{quando} ela é necessária frente a logs assinados ou
        bancos append-only, e a delimitação do que o protótipo atual
        ainda não implementa; e
  \item uma avaliação \emph{offline} reprodutível (código, dados,
        modelos e hiperparâmetros publicados) do motor de risco sobre
        dados sintéticos.
\end{itemize}

O restante do artigo organiza-se assim. A Seção~\ref{sec:related} revisa
os trabalhos correlatos. A Seção~\ref{sec:arch} descreve a arquitetura
proposta; a Seção~\ref{sec:ai}, o motor de IA; e a Seção~\ref{sec:bc}, a
camada de auditoria em blockchain. A Seção~\ref{sec:method} apresenta a
metodologia de avaliação e a Seção~\ref{sec:results}, os resultados. A
Seção~\ref{sec:conclusion} conclui.

\section{Trabalhos Correlatos}
\label{sec:related}
\textbf{Biometria comportamental.} A autenticação contínua a partir de
padrões de interação está bem estabelecida. O TypeNet~\cite{acien2022} e
o TypeFormer~\cite{stragapede2024} identificam usuários pela dinâmica de
digitação com baixas taxas de erro, e o Touchalytics~\cite{frank2013}
mostra que padrões de toque discriminam usuários em dispositivos móveis.
Esses trabalhos confirmam o poder discriminativo do comportamento do
usuário, mas concentram-se em uma única modalidade biométrica e não
tratam da auditoria imutável das decisões. Surveys recentes de
autenticação contínua baseada em sensores~\cite{abuhamad2021} corroboram
que sinais comportamentais generalizam por diversas modalidades, ao
mesmo tempo em que apontam desafios em aberto de robustez, envelhecimento
de template e implantação real --- desafios que reforçam a fusão do
comportamento com outros sinais contextuais em vez da dependência de uma
única modalidade.

\textbf{Sinais contextuais e de dispositivo.} Além do comportamento, o
contexto de dispositivo e rede de um login é, por si só, discriminativo.
Alaca e van Oorschot~\cite{alaca2016} classificam vinte e nove métodos de
fingerprinting de dispositivo para reforçar a autenticação web,
comparando-os por estabilidade, dificuldade de spoofing e
distinguibilidade, e Andriamilanto \emph{et al.}~\cite{andriamilanto2021}
mostram, em mais de quatro milhões de navegadores, que fingerprints de
navegador carregam distinguibilidade e estabilidade suficientes para
reforçar a autenticação. Freeman \emph{et al.}~\cite{freeman2016}
introduziram o primeiro modelo probabilístico público que classifica um
login como normal ou suspeito a partir de sinais como IP de origem,
geolocalização, configuração do navegador e hora do dia. Esses trabalhos
validam sinais contextuais individuais, mas os tratam em grande parte de
forma isolada e não acoplam a decisão resultante a um registro à prova de
adulteração.

\textbf{Autenticação baseada em risco.} Wiefling \emph{et
al.}~\cite{wiefling2019} conduziram o primeiro estudo empírico de larga
escala sobre RBA em serviços reais, e trabalhos posteriores
\cite{wiefling2023} mostraram que os fatores de risco devem ser
cuidadosamente ajustados por serviço. Nossa proposta segue essa linha,
mas substitui heurísticas fixas --- e os modelos de sinal único,
pontuados estatisticamente, que as precederam~\cite{freeman2016} --- por
um ensemble de ML multissinal e adiciona uma camada de não repúdio.

\textbf{Detecção de anomalias.} O Isolation Forest~\cite{liu2008} é
referência para detecção não supervisionada de outliers em alta
dimensionalidade, e modelos supervisionados e profundos
\cite{goodfellow2016} o complementam quando há rótulos; Killourhy e
Maxion~\cite{killourhy2009} estabeleceram um marco metodológico para a
detecção de anomalias em dinâmica de digitação. Combinamos ambas as
naturezas --- não supervisionada e supervisionada --- em um ensemble.

\textbf{Blockchain para auditoria.} Embora a auditoria de logs baseada em
blockchain venha sendo investigada, a maioria das abordagens trata o
registro de eventos de forma \emph{desacoplada} do mecanismo de decisão.
Sistemas como o BlockAudit~\cite{blockaudit} demonstram trilhas de
auditoria escaláveis e à prova de adulteração em blockchain, mas os
eventos registrados originam-se em aplicações externas à camada de
registro, de modo que o registro atesta \emph{que} um evento ocorreu sem
vinculá-lo ao raciocínio que o produziu. O diferencial pretendido por este
trabalho é integrar, em um único fluxo, a decisão adaptativa por IA e o seu registro
imutável em rede permissionada, de modo que cada entrada armazenada
carregue os próprios critérios e o escore que justificaram a ação. A
Tabela~\ref{tab:related} sintetiza trabalhos representativos: nenhuma
linha de pesquisa reúne simultaneamente (i)~a fusão de múltiplos sinais
contextuais e comportamentais, (ii)~a decisão adaptativa por ensemble e
(iii)~a auditoria imutável acoplada à decisão.

\begin{table}[t]
\caption{Posicionamento Frente a Trabalhos Correlatos Representativos}
\label{tab:related}
\centering
\footnotesize
\setlength{\tabcolsep}{3.5pt}
\renewcommand{\arraystretch}{1.2}
\begin{tabular}{@{}p{1.85cm}p{1.3cm}ccc@{}}
\toprule
\textbf{Trabalho} & \textbf{Sinais} & \textbf{Adapt.} & \textbf{Audit.} & \textbf{Dados} \\
\midrule
TypeNet~\cite{acien2022} & Digitação & \checkmark & --- & Real \\
Touchalytics~\cite{frank2013} & Toque & \checkmark & --- & Real \\
Freeman~\cite{freeman2016} & IP/geo/UA & Parc. & --- & Real \\
Alaca~\cite{alaca2016} & Fp.\ disp. & --- & --- & --- \\
Andriam.~\cite{andriamilanto2021} & Fp.\ nav. & --- & --- & Real \\
RBA~\cite{wiefling2019} & IP/contexto & Parc. & --- & Real \\
Iso.\ Forest~\cite{liu2008} & Genérico & N.\ sup. & --- & --- \\
BlockAudit~\cite{blockaudit} & Logs & --- & \checkmark & --- \\
\textbf{Este trab.} & \textbf{14 ctx.+comp.} & \textbf{Ens.} & \textbf{Proj.} & Sint. \\
\bottomrule
\end{tabular}

\vspace{2pt}
{\footnotesize \emph{Adapt.}: decisão adaptativa/aprendida (Parc.:
parcial; N.\ sup.: não supervisionada; Ens.: ensemble). \emph{Audit.}:
auditoria imutável acoplada à decisão (Proj.: projetada, ainda não
implementada).}
\end{table}

\section{Arquitetura Proposta}
\label{sec:arch}
A arquitetura articula microsserviços, gestão federada de identidade,
análise de risco por IA e registro imutável em blockchain, seguindo os
princípios de segurança por desenho e verificação contínua do Zero
Trust~\cite{rose2020,nist80053r5}. Em vez de uma decisão binária no
momento do login, o modelo adota uma perspectiva contínua e orientada a
risco. Organiza-se em cinco camadas integradas (Fig.~\ref{fig:arch}):

\begin{enumerate}
  \item \textbf{Aplicação Modular.} Serviços independentes (Spring Boot)
        e front end Angular expondo APIs REST autenticadas.
  \item \textbf{Identidade e Acesso.} Autenticação federada via Keycloak
        (OAuth2 / OpenID Connect), emissão de JWT e suporte a MFA.
  \item \textbf{Inteligência Contextual.} Extração dos quatorze critérios
        comportamentais e cálculo do escore de risco pelo ensemble.
  \item \textbf{Decisão Adaptativa.} Aplicação dos limiares de risco,
        resultando em \emph{permitir}, \emph{reforçar com MFA} ou
        \emph{bloquear}.
  \item \textbf{Registro Imutável e Auditoria.} Persistência de cada
        decisão em blockchain permissionada, garantindo rastreabilidade
        e não repúdio.
\end{enumerate}

O fluxo operacional é uma cadeia decisória encadeada: a tentativa de
login dispara a extração dos critérios contextuais; o motor de IA calcula
o escore de risco; os limiares de decisão são aplicados; o resultado é
registrado na blockchain; e o sistema permite o acesso, exige MFA ou o
bloqueia. A separação em camadas permite evoluir cada responsabilidade de
forma independente e aplicar políticas de segurança específicas em cada
nível.

\textbf{Modelo de ameaças e racional de projeto.} O sistema pressupõe um
adversário que pode já deter credenciais válidas --- obtidas por
vazamento, phishing ou reutilização --- e que pode tentar logins a partir
de redes e dispositivos arbitrários, mas que não consegue forjar o
consenso da rede permissionada nem observar o estado privado dos modelos
de risco. Sob esse modelo, as camadas tratam preocupações complementares:
a C2 estabelece \emph{quem} alega autenticar-se, a C3 estima \emph{quão
provável} é que a alegação seja legítima dado o contexto, a C4 converte a
estimativa em uma ação proporcional e a C5 torna a ação
\emph{prestável de contas}. A faixa intermediária pressupõe que o
adversário não detém o segundo fator; um adversário com credenciais e
acesso ao segundo fator (p.\,ex., por phishing em tempo real) só é
contido pela faixa de bloqueio, razão pela qual a proporção de ataques
que cai em cada faixa precisa ser reportada, e não apenas o desempenho
binário. Desacoplar a estimativa (C3) da aplicação
(C4) é deliberado: os limiares de risco podem ser recalibrados, e os
modelos reentreinados, sem tocar no provedor de identidade, e a camada de
auditoria permanece agnóstica aos aprendizes específicos em uso. Como
cada decisão é registrada junto com sua justificativa --- o hash dos
critérios de entrada e o escore que o produziu ---, um auditor pode
reproduzir posteriormente qualquer decisão e verificar que a ação
registrada corresponde à política vigente naquele momento, sem precisar
confiar nos próprios logs do operador. É esse acoplamento estreito entre
\emph{decidir} e \emph{registrar} que distingue a proposta de pipelines
que registram eventos depois, e separadamente, da decisão.

\textbf{Estado do protótipo.} A arquitetura acima é o alvo do projeto. O
protótipo público (Spring Boot + Angular) implementa até agora: login com
senha e JWT próprios (a integração com o Keycloak existe como perfil de
configuração opcional, não utilizado na implantação atual); MFA por TOTP;
uma versão em Java dos três modelos, treinada na inicialização; e a
persistência de cada decisão em uma tabela relacional (PostgreSQL) com um
encadeamento de hashes. Três divergências em relação ao texto são
explicitadas: (i)~os modelos avaliados na Seção~\ref{sec:results} são os
do pipeline \emph{offline} em Python, e não os do protótipo Java, que usa
outro conjunto de quatorze atributos; (ii)~no protótipo, o escore do
ensemble ainda é combinado a um escore de confiança
($T = 0{,}4\,H + 0{,}3\,(1-A) + 0{,}2\,R + 0{,}1\,X$, com $H$ histórico
de sucesso, $A$ risco da IA, $R$ comportamento recente e $X$ fatores
externos), e o MFA é exigido quando $T<0{,}5$, sem faixa de bloqueio;
(iii)~a camada C5 não está conectada a uma rede Fabric: o serviço
correspondente é um simulador, e a verificação de integridade da cadeia
de hashes ainda não é funcional. Alinhar o protótipo à
Eq.~(\ref{eq:ensemble}) e às três faixas é o próximo passo de
implementação.

\begin{figure}[t]
\centering
\begin{tikzpicture}[
  layer/.style={draw, rounded corners, minimum width=7.4cm,
    minimum height=0.72cm, align=center, font=\footnotesize, fill=gray!8},
  arr/.style={-{Latex[length=2mm]}, thick}
]
\node[layer] (l1) {C1 -- Aplicação Modular \\ \scriptsize Serviços Spring Boot $\cdot$ front end Angular $\cdot$ APIs REST};
\node[layer, below=0.28cm of l1] (l2) {C2 -- Identidade e Acesso \\ \scriptsize Keycloak $\cdot$ OAuth2 / OIDC $\cdot$ JWT $\cdot$ MFA};
\node[layer, below=0.28cm of l2] (l3) {C3 -- Inteligência Contextual \\ \scriptsize 14 critérios $\cdot$ ensemble IF + RF + DL $\rightarrow$ escore};
\node[layer, below=0.28cm of l3] (l4) {C4 -- Decisão Adaptativa \\ \scriptsize limiares $\rightarrow$ permitir / MFA / bloquear};
\node[layer, below=0.28cm of l4] (l5) {C5 -- Registro Imutável e Auditoria \\ \scriptsize Hyperledger Fabric $\cdot$ trilha de auditoria verificável};
\draw[arr] (l1) -- (l2);
\draw[arr] (l2) -- (l3);
\draw[arr] (l3) -- (l4);
\draw[arr] (l4) -- (l5);
\end{tikzpicture}
\caption{Arquitetura em cinco camadas do sistema proposto.}
\label{fig:arch}
\end{figure}

\section{Motor de Decisão por IA}
\label{sec:ai}
\textbf{Ensemble.} O motor de risco combina três aprendizes
complementares. O \emph{Isolation Forest}~\cite{liu2008} atua como
primeiro nível de análise, isolando outliers no espaço multidimensional
de critérios com poucas partições aleatórias; contribui com 40\% do peso
final. O \emph{Random Forest} contribui com 30\%, capturando relações não
lineares entre os critérios e mantendo-se interpretável. Uma \emph{rede
neural multicamada}~\cite{goodfellow2016} contribui com os 30\%
restantes, com camada de entrada de 14 neurônios, duas camadas ocultas
(64 e 32 unidades, ativações ReLU, otimizador Adam, sem dropout) e saída
logística. O escore final é a média ponderada
\begin{equation}
S = 0{,}4\,S_{\mathrm{IF}} + 0{,}3\,S_{\mathrm{RF}} + 0{,}3\,S_{\mathrm{DL}},
\label{eq:ensemble}
\end{equation}
em que cada componente retorna um valor em $[0,1]$ (0: normal, 1:
altamente anômalo). A combinação de aprendizes não supervisionados e
supervisionados visa cobrir tanto desvios inéditos (sem rótulo) quanto
padrões de fraude conhecidos; a contribuição de cada modelo é medida na
Seção~\ref{sec:results}.

\textbf{Critérios.} Cada tentativa é representada por quatorze critérios
em quatro famílias: \emph{temporais} (hora do acesso,
dia da semana, intervalo desde o último login, frequência em 24\,h);
\emph{geográficos/rede} (distância geodésica do local habitual, tipo de
rede, ASN, velocidade de deslocamento impossível);
\emph{dispositivo/navegador} (fingerprint do navegador, sistema
operacional, tipo de dispositivo, detecção de automação/headless); e
\emph{comportamento de sessão} (padrão de digitação e histórico de falhas
recentes). Esses sinais visam credential stuffing, viagem impossível,
sequestro de sessão, bots e desvios zero-day.

\textbf{Limiares adaptativos.} O escore é mapeado em três faixas de
risco, cujos cortes (0,3 e 0,7) são parâmetros de projeto, e não valores
calibrados:
baixo risco ($S<0{,}3$) permite o acesso sem etapas adicionais; risco
moderado ($0{,}3\le S<0{,}7$) aciona MFA; e alto risco ($S\ge0{,}7$)
bloqueia a tentativa, registra o evento na blockchain e notifica o
administrador de segurança. A avaliação da
Seção~\ref{sec:results} usa, além das faixas, um limiar binário único
($\tau=0{,}54$, Seção~\ref{sec:method}), que cai dentro da faixa de MFA;
por isso precisão, recall e FPR não descrevem, sozinhos, a liberação nem
o bloqueio, e o desempenho por faixa é reportado à parte
(Tabela~\ref{tab:bands}).
Os três modelos são avaliados sequencialmente. Retreinamento periódico
frente à deriva de conceito~\cite{gama2014} e versionamento de modelos
são requisitos de projeto ainda não implementados.

\section{Auditoria Baseada em Blockchain}
\label{sec:bc}
Esta seção descreve o \emph{projeto} da camada de auditoria; como
explicado na Seção~\ref{sec:arch}, ela ainda não foi integrada a uma rede
real, e sua avaliação é trabalho futuro. No projeto, cada decisão é
registrada em uma rede permissionada
baseada em Hyperledger Fabric~\cite{fabric2024}, cujo controle de
identidade e governança estruturada são adequados a ambientes
corporativos e governamentais~\cite{nistir8202}. Para cada evento, o
registro armazena o hash dos critérios de entrada, o escore de risco, a
decisão (permitir, MFA ou bloquear) e o timestamp da transação, encadeado
ao registro anterior de modo a formar uma cadeia cronológica de
evidências.

A escolha do blockchain frente a mecanismos mais simples é justificada, e
não presumida. Logs assinados digitalmente e bancos append-only já
oferecem integridade e não repúdio por meio de assinaturas
criptográficas, sendo suficientes quando existe uma autoridade central
confiável que opera o registro. A diferença essencial de um blockchain
permissionado \emph{não} está na assinatura em si, mas em não depender da
confiança em um único custodiante: o registro é replicado e validado por
consenso entre múltiplos participantes, tornando-se resistente à
adulteração mesmo por administradores privilegiados do próprio sistema e
permitindo verificação independente por auditores externos. Isso é
determinante aqui, pois as decisões de acesso são tomadas de forma
autônoma por modelos de IA e precisam ser auditáveis por terceiros (por
exemplo, em fiscalizações de conformidade da LGPD~\cite{lgpd2018}),
cenário em que o operador, sozinho, não deve poder alterar a trilha de
decisões. Reconhece-se o custo --- maior complexidade operacional e menor
throughput do que um banco append-only ---, que se justifica justamente
em cenários multipartes e regulados, delimitando o escopo em que o
blockchain é efetivamente necessário. Esse custo ainda não foi medido; a
avaliação planejada exige uma rede com ao menos duas organizações e
mede a latência de \emph{commit}, a vazão sob carga e a detecção de uma
adulteração deliberada por um administrador de uma das organizações.
Enquanto isso não for feito, o registro relacional do protótipo continua
sob controle de um único custodiante --- exatamente o cenário que esta
seção considera insuficiente.

\section{Metodologia de Avaliação}
\label{sec:method}
\textbf{Dados e rotulagem.} Diante da ausência de bases públicas
rotuladas para autenticação contextual, construiu-se um conjunto de dados
sintéticos com 50.000 registros de login,\footnote{Conjunto de dados (CSV rotulado, modelos treinados e métricas): \url{https://github.com/jgbcientista/gestao-autonomo/tree/main/docs/artefatos-modelos}} com
proporção de 95\% de acessos legítimos e 5\% suspeitos, fixada
como hipótese de trabalho (não se trata de uma distribuição medida em
produção). Cada registro contém os quatorze critérios da
Seção~\ref{sec:ai}. \emph{O rótulo de cada registro é o ramo do gerador
que o produziu} (legítimo ou perturbado); não houve revisão manual por
especialistas. Todas as distribuições são constantes escolhidas pelos
autores para representar um uso corporativo plausível, e não parâmetros
estimados a partir de dados reais: a hora de acesso é uma gaussiana
($\mu=13$\,h, $\sigma=3{,}2$\,h) truncada em $[0,23]$ --- tratada como
variável linear, sem codificação circular ---, o intervalo entre logins
segue uma exponencial, a frequência em 24\,h uma Poisson e a distância
geodésica uma meia-normal; os escores de fingerprint do navegador e de
similaridade de digitação são gaussianas de média alta truncadas em
$[0,1]$, e os indicadores de dispositivo, sistema operacional, ASN e
tipo de rede são variáveis categóricas com frequências fixadas no
gerador. Os registros suspeitos são obtidos perturbando \emph{de dois a
quatro} critérios, escolhidos ao acaso, de um perfil legítimo para faixas
anômalas (horário de madrugada, fim de semana, rajada de logins,
distância e velocidade de deslocamento impossíveis, rede ou ASN
atípicos, sistema operacional ou dispositivo desconhecidos,
fingerprint e digitação divergentes, automação, falhas recentes). Cerca de 5\% dos registros legítimos recebem uma perturbação
em \emph{um} critério, sorteada \emph{das mesmas faixas} usadas para os
ataques; é essa escolha que gera os falsos positivos. Por fim,
adiciona-se a todos os critérios ruído gaussiano com desvio de $0{,}30$
vezes o desvio-padrão da coluna. Como o ruído é aplicado também às
variáveis categóricas e às não negativas, o conjunto contém valores
fisicamente implausíveis (p.\,ex., ASN $=0{,}94$ ou velocidade de
deslocamento negativa); registramos essa limitação do gerador em vez de
ocultá-la, e sua correção faz parte dos trabalhos futuros. Por se tratar de dados
sintéticos, não há coleta de
informações pessoais identificáveis nesta fase; a evolução para dados
reais seguirá os princípios de minimização e anonimização exigidos pela
LGPD~\cite{lgpd2018}.

\textbf{Protocolo.} O conjunto foi particionado de forma estratificada em
70\% para treinamento e 30\% para teste. O normalizador (padronização) e
a normalização min-max do escore da Isolation Forest são ajustados
apenas no treino. Os hiperparâmetros (Tabela~\ref{tab:hyper}) foram
fixados \emph{a priori}, sem busca; o único parâmetro ajustado é o limiar
de decisão, escolhido na grade $\{0{,}20, 0{,}21, \dots, 0{,}80\}$ como o
que maximiza o F1 no treino ($\tau=0{,}54$) e aplicado sem alteração ao
teste. Para estimar a variabilidade, executou-se uma validação cruzada
estratificada de 5 folds \emph{apenas sobre a partição de treino}, em que
o normalizador, os modelos e o limiar são reajustados dentro de cada
fold; a partição de teste não participa. O desempenho foi medido por
precisão, recall, F1-score, taxa de falsos positivos (FPR), AUC-ROC,
AUC da curva precisão-recall (PR-AUC, mais informativa sob desbalanceamento)
e latência de inferência. Para a decisão em três faixas, definem-se a
\emph{taxa de atrito} (fração de acessos legítimos que recebem MFA ou
bloqueio) e a \emph{taxa de passagem silenciosa} (fração de ataques
liberados sem nenhum desafio). As sementes aleatórias foram fixadas.

\textbf{Análises adicionais.} Com o mesmo pipeline, avaliam-se ainda:
(i)~cada modelo isolado, com limiar próprio escolhido no treino;
(ii)~as 66 combinações de pesos em passos de 0,1 (análise de
sensibilidade, não de seleção); (iii)~proporções de ataques de 1\%, 2\%,
5\%, 10\% e 20\%; (iv)~uma variante do gerador sem ruído nas variáveis
categóricas e de contagem e com truncamento das contínuas não negativas,
que elimina os valores implausíveis; e (v)~um atacante que conhece os
quatorze critérios e substitui os que controla por valores sorteados da
distribuição legítima, em quatro níveis cumulativos: horário e ritmo de
tentativas; rede, ASN e localização (proxy residencial próximo à vítima);
fingerprint, sistema, dispositivo e automação (navegador
anti-detecção); e padrão de digitação.
Para permitir verificação independente, o código-fonte,
os artefatos dos modelos treinados e os scripts que geram o conjunto de
dados e reproduzem cada experimento estão abertamente
disponíveis.\footnote{\url{https://github.com/jgbcientista/gestao-autonomo}}

\textbf{Reprodutibilidade.} Um único script de entrada regenera todo o
pipeline de forma determinística a partir da semente fixa: o conjunto
rotulado (50.000 registros com os quatorze critérios e o rótulo de
classe), os três modelos treinados juntamente com o normalizador de
atributos (serializados com \texttt{joblib}), a configuração do ensemble
(pesos dos componentes, os limites de normalização do escore da
Isolation Forest e o limiar de decisão selecionado pelo melhor F1 de
treinamento) e um arquivo de métricas legível por máquina contendo cada
valor reportado na Seção~\ref{sec:results}. Reexecutar o script em um
ambiente limpo reproduz as métricas reportadas; diferenças de versão do
scikit-learn produzem apenas variações na terceira casa decimal (p.\,ex.,
F1 de 0,802 a 0,805), de
modo que um revisor possa auditar todo o pipeline de ponta a ponta ---da
geração dos atributos brutos ao limiar de decisão final--- sem acesso a
quaisquer dados privados ou pessoalmente identificáveis. Esse
empacotamento é o que transforma a descrição de ``cerca de 50.000
registros'' em um artefato que um terceiro pode regenerar, inspecionar e
estender.

\begin{table}[t]
\caption{Hiperparâmetros Utilizados (fixados a priori)}
\label{tab:hyper}
\centering
\footnotesize
\setlength{\tabcolsep}{3pt}
\renewcommand{\arraystretch}{1.15}
\begin{tabular}{@{}p{1.9cm}p{6.0cm}@{}}
\toprule
\textbf{Componente} & \textbf{Configuração} \\
\midrule
Pré-processamento & padronização (z-score) ajustada no treino \\
Isolation Forest & 100 árvores; contaminação 0,05; escore $-$\texttt{score\_samples} normalizado min-max pelo treino e truncado em $[0,1]$ \\
Random Forest & 120 árvores; profundidade livre; \texttt{class\_weight=balanced} \\
Rede neural (MLP) & camadas ocultas 64--32, ReLU, Adam, 220 iterações, sem dropout; saída logística \\
Ensemble & $0{,}4\,S_{\mathrm{IF}}+0{,}3\,S_{\mathrm{RF}}+0{,}3\,S_{\mathrm{DL}}$ \\
Limiar & grade 0,20--0,80 (passo 0,01), máx.\ F1 no treino $\Rightarrow\tau=0{,}54$ \\
Dados & 50.000 registros; 5\% suspeitos; ruído $0{,}30\sigma$; semente 42 \\
\bottomrule
\end{tabular}
\end{table}

\textbf{Baselines.} A Tabela~\ref{tab:baseline} confronta a proposta com
baselines representativos da literatura e da prática. Os valores dos
baselines provêm da literatura citada e não foram reexecutados aqui; a
comparação é, portanto, qualitativa quanto às capacidades de cada
abordagem. Os três eixos de comparação não são arbitrários; decorrem
dos requisitos que a literatura impõe à autenticação moderna. A
\emph{sensibilidade ao contexto} (Ctx.) reflete a recomendação do NIST
e do modelo Zero Trust de que a verificação seja condicionada aos sinais
de cada tentativa, e não apenas à credencial
\cite{nist80063b,rose2020}, propriedade que estudos empíricos de
autenticação baseada em risco apontam como decisiva na
prática~\cite{wiefling2019,freeman2016}. A \emph{adaptabilidade}
(Adapt.) expressa se a resposta é proporcional ao risco estimado ---
princípio central da autenticação baseada em
risco~\cite{nist80063b,wiefling2019,wiefling2023} --- em vez de uma
regra fixa. A \emph{auditabilidade} (Audit.) trata da prestação de
contas de decisões automatizadas sobre dados pessoais, exigência
explicitada por regulações como a LGPD~\cite{lgpd2018} e, de modo mais
amplo, pela necessidade de verificação independente quando o acesso é
delegado a modelos autônomos. Em conjunto, esses eixos cobrem as duas
propriedades que um sistema de autenticação contextual precisa
satisfazer simultaneamente --- decidir corretamente e ser verificável
---, enquanto a coluna \emph{principal limitação} indica onde cada
baseline falha em ao menos uma delas.

\begin{table}[t]
\caption{Comparação com Abordagens de Referência}
\label{tab:baseline}
\centering
\footnotesize
\setlength{\tabcolsep}{3.5pt}
\renewcommand{\arraystretch}{1.2}
\begin{tabular}{@{}p{1.7cm}cccp{2.0cm}@{}}
\toprule
\textbf{Abordagem} & \textbf{Ctx.} & \textbf{Adapt.} & \textbf{Audit.} & \textbf{Limitação principal} \\
\midrule
Senha + 2FA estático & --- & --- & --- & Ignora contexto; phishing/SIM swap \\
RBA por IP~\cite{wiefling2019} & Parc. & Heur. & --- & Feature única; ajuste manual \\
Escore estatístico~\cite{freeman2016} & \checkmark & Est. & --- & Poucos sinais; sem auditoria \\
Keystroke isolado~\cite{killourhy2009} & \checkmark & --- & --- & Offline, fator único (EER $\approx$9,6\%) \\
\textbf{Proposta} & \textbf{\checkmark} & \textbf{Ens.} & \textbf{Proj.} & Dados sintéticos; auditoria em projeto \\
\bottomrule
\end{tabular}

\vspace{2pt}
{\footnotesize \emph{Ctx.}: usa contexto do login; \emph{Adapt.}: decisão
adaptativa (Heur.: heurística; Est.: estatística; Ens.: ensemble);
\emph{Audit.}: auditoria imutável acoplada à decisão (Proj.: projetada).}
\end{table}

\section{Resultados e Discussão}
\label{sec:results}
Sobre o conjunto de teste (30\% dos dados, não vistos durante o
treinamento), o ensemble alcançou acurácia de 98,1\%, precisão de 83,0\%
e recall de 77,6\% na detecção de acessos suspeitos, com F1-score de
80,2\% e AUC-ROC de 0,99, indicando forte capacidade de ordenação entre
acessos legítimos e suspeitos. A taxa de falsos positivos manteve-se
baixa, em 0,8\%, fator relevante para a usabilidade. A validação cruzada
restrita ao treino confirma a ordem de grandeza (F1 $80,3\%\pm1,3\%$;
AUC $0,990\pm0,002$), com limiares entre 0,42 e 0,56 conforme o fold. A
matriz de confusão no teste é TN\,=\,14.131, FP\,=\,119, FN\,=\,168,
TP\,=\,582. A Tabela~\ref{tab:results} resume esses números.

\begin{table}[t]
\caption{Desempenho do Ensemble no Conjunto de Teste}
\label{tab:results}
\centering
\renewcommand{\arraystretch}{1.2}
\begin{tabular}{@{}lc@{\hspace{5mm}}lc@{}}
\toprule
\textbf{Métrica} & \textbf{Valor} & \textbf{Métrica} & \textbf{Valor} \\
\midrule
Acurácia   & 98,1\% & F1-score        & 80,2\% \\
Precisão   & 83,0\% & AUC-ROC         & 0,99   \\
Recall     & 77,6\% & Taxa falsos pos.\ & 0,8\% \\
\midrule
\multicolumn{2}{@{}l}{F1 (CV 5 folds no treino)} & \multicolumn{2}{l}{$80,3\%\pm1,3\%$} \\
\multicolumn{2}{@{}l}{AUC (CV 5 folds no treino)} & \multicolumn{2}{l}{$0,990\pm0,002$} \\
\multicolumn{2}{@{}l}{Latência do ensemble (\emph{offline})} & \multicolumn{2}{l}{$\approx$33\,ms/req.} \\
\bottomrule
\end{tabular}
\end{table}

\textbf{Decisão em três faixas.} A Tabela~\ref{tab:bands} aplica as faixas
0,3/0,7 ao escore do teste. Entre os acessos legítimos, 97,3\% são
liberados sem desafio (taxa de atrito de 2,7\%) e apenas 0,1\% são
bloqueados. Entre os ataques, 62,5\% são bloqueados, 30,1\% recebem MFA e
7,3\% passam silenciosamente. Como a faixa de MFA só contém um adversário
que não detém o segundo fator (Seção~\ref{sec:arch}), a proteção efetiva
contra quem tem credenciais e segundo fator é a fração bloqueada, e não o
recall binário.

\begin{table}[t]
\caption{Decisão em Três Faixas no Conjunto de Teste}
\label{tab:bands}
\centering
\footnotesize
\renewcommand{\arraystretch}{1.15}
\begin{tabular}{@{}lccc@{}}
\toprule
\textbf{Classe real} & \textbf{Permitir} & \textbf{MFA} & \textbf{Bloquear} \\
 & ($S<0{,}3$) & & ($S\ge0{,}7$) \\
\midrule
Legítimo & 97,3\% & 2,6\% & 0,1\% \\
Suspeito & 7,3\% & 30,1\% & 62,5\% \\
\midrule
\multicolumn{4}{@{}l}{Taxa de atrito: 2,7\% \quad Taxa de passagem silenciosa: 7,3\%} \\
\bottomrule
\end{tabular}
\end{table}

\textbf{Contribuição de cada modelo.} A Tabela~\ref{tab:single} compara o
ensemble a cada modelo treinado isoladamente no mesmo pipeline. O ensemble
ganha cerca de 3 pontos de F1 sobre o melhor modelo isolado (0,802 contra
0,775 da Random Forest), reduz a FPR de 1,15\% para 0,84\% e eleva a
PR-AUC de 0,86 para 0,90. Os pesos 0,4/0,3/0,3, porém, não foram
otimizados: nas 66 combinações avaliadas, o F1 no teste varia de 0,767
(só Isolation Forest) a 0,820, e a configuração adotada fica em 19.º
lugar. As melhores combinações atribuem 0,6--0,8 à Isolation Forest e
pouco ou nenhum peso à rede neural. Escolher os pesos por esse resultado
seria ajustar ao teste; a seleção correta, em uma partição de validação,
fica para a próxima etapa.

\begin{table}[t]
\caption{Ensemble versus Modelos Isolados (Teste)}
\label{tab:single}
\centering
\footnotesize
\setlength{\tabcolsep}{4pt}
\renewcommand{\arraystretch}{1.15}
\begin{tabular}{@{}lcccccc@{}}
\toprule
\textbf{Modelo} & $\tau$ & \textbf{Prec.} & \textbf{Rec.} & \textbf{F1} & \textbf{FPR} & \textbf{PR-AUC} \\
\midrule
Isolation Forest & 0,45 & 0,763 & 0,772 & 0,767 & 1,26\% & 0,855 \\
Random Forest    & 0,60 & 0,779 & 0,771 & 0,775 & 1,15\% & 0,857 \\
Rede neural      & 0,48 & 0,757 & 0,792 & 0,774 & 1,34\% & 0,860 \\
\textbf{Ensemble} & 0,54 & \textbf{0,830} & 0,776 & \textbf{0,802} & \textbf{0,84\%} & \textbf{0,903} \\
\bottomrule
\end{tabular}
\end{table}

\textbf{Sensibilidade à proporção de ataques e ao gerador.} A
Tabela~\ref{tab:prop} repete o experimento para outras proporções. A AUC
permanece em 0,99, mas as métricas dependentes da prevalência mudam
bastante: com 1\% de ataques, o F1 cai para 0,63 e a passagem silenciosa
sobe para 24,0\%, porque a Isolation Forest (contaminação igual à
proporção) e o limiar escolhido por F1 passam a favorecer a classe
majoritária. Em produção, em que ataques tendem a ser mais raros que 5\%,
esse é o cenário relevante. Já a correção do ruído do gerador (item~(iv)
das análises adicionais) quase não altera os resultados (F1 0,809, AUC
0,992, passagem silenciosa 6,9\%), o que indica que os valores
implausíveis não inflavam as métricas.

\begin{table}[t]
\caption{Sensibilidade à Proporção de Ataques}
\label{tab:prop}
\centering
\footnotesize
\setlength{\tabcolsep}{4pt}
\renewcommand{\arraystretch}{1.15}
\begin{tabular}{@{}lccccc@{}}
\toprule
\textbf{Ataques} & \textbf{F1} & \textbf{AUC} & \textbf{PR-AUC} & \textbf{Atrito} & \textbf{Pass.\ silenc.} \\
\midrule
1\%  & 0,632 & 0,989 & 0,718 & 0,8\% & 24,0\% \\
2\%  & 0,720 & 0,991 & 0,809 & 1,1\% & 16,7\% \\
5\%  & 0,802 & 0,992 & 0,903 & 2,7\% & 7,3\% \\
10\% & 0,869 & 0,992 & 0,944 & 3,1\% & 6,4\% \\
20\% & 0,917 & 0,993 & 0,976 & 4,2\% & 3,7\% \\
\bottomrule
\end{tabular}
\end{table}

\textbf{Evasão adversarial.} A Tabela~\ref{tab:evasion} mostra o efeito de
um atacante que conhece os critérios. Imitar apenas o horário e o ritmo
das tentativas --- algo trivial --- já reduz o recall de 77,6\% para
50,7\% e eleva a passagem silenciosa de 7,3\% para 28,4\%. Com um proxy
residencial próximo à vítima, o recall cai para 18,0\% e 64,4\% dos
ataques passam sem desafio; com um navegador anti-detecção, praticamente
todos passam. Esse resultado é, em parte, consequência do gerador: os
ataques sintéticos \emph{são definidos} como desvios nesses critérios, de
modo que removê-los apaga o sinal. Mas ele também expõe um limite real
da abordagem: a maioria dos quatorze critérios é controlável pelo
adversário, e só o padrão de digitação e o histórico da conta resistem a
um atacante informado. Critérios difíceis de forjar e defesas contra
evasão são, portanto, requisito, e não refinamento.

\begin{table}[t]
\caption{Evasão por Atacante que Conhece os Critérios (Ataques do Teste)}
\label{tab:evasion}
\centering
\footnotesize
\setlength{\tabcolsep}{3pt}
\renewcommand{\arraystretch}{1.15}
\begin{tabular}{@{}p{3.3cm}cccc@{}}
\toprule
\textbf{Critérios imitados (cumulativo)} & \textbf{Recall} & \textbf{Permitir} & \textbf{MFA} & \textbf{Bloquear} \\
\midrule
Nenhum & 77,6\% & 7,3\% & 30,1\% & 62,5\% \\
+ horário e ritmo & 50,7\% & 28,4\% & 33,1\% & 38,5\% \\
+ rede, ASN, local (proxy) & 18,0\% & 64,4\% & 25,3\% & 10,3\% \\
+ fingerprint, SO, disp., automação & 2,9\% & 91,1\% & 8,0\% & 0,9\% \\
+ digitação (todos) & 0,3\% & 98,1\% & 1,9\% & 0,0\% \\
\bottomrule
\end{tabular}

\vspace{2pt}
{\footnotesize Recall no limiar binário $\tau=0{,}54$; demais colunas com as faixas 0,3/0,7.}
\end{table}

A latência foi medida apenas no pipeline \emph{offline}
(scikit-learn, uma amostra por chamada, mediana de 300 chamadas): 4,0\,ms
para a Isolation Forest, 29,2\,ms para a Random Forest e 0,05\,ms para a
rede neural, executadas em sequência, totalizando cerca de 33\,ms --- valor
dependente da máquina e dominado pela Random Forest. Não foram medidos a
latência do fluxo completo de login, o custo do registro de auditoria nem
a vazão; esses números dependem da integração descrita na
Seção~\ref{sec:bc} e ficam para a próxima etapa.

\textbf{Discussão.} Os resultados parciais são promissores e sugerem a
viabilidade técnica da abordagem; contudo, por decorrerem de dados
sintéticos e de uma fase preliminar, não permitem ainda confirmar a
hipótese central de forma conclusiva. Uma comparação rigorosa com os
baselines e a validação com dados reais permanecem necessárias antes de
qualquer afirmação quantitativa de superioridade. Ainda assim, em
ambiente controlado, o ensemble separou os registros perturbados dos
legítimos com AUC de 0,99, o que indica que os critérios escolhidos
carregam sinal quando os ataques se manifestam nas faixas modeladas pelo
gerador --- ressalva que a análise de evasão torna concreta. Nenhuma afirmação sobre
resistência à adulteração é feita nesta versão, pois a camada de
auditoria não foi avaliada.

\textbf{Análise de erros e compromissos.} O limiar de decisão não é
fixado a priori, mas selecionado como o valor que maximiza o F1-score na
partição de treinamento, o que torna o ponto de operação explícito e
ajustável por implantação: reduzi-lo aumenta o recall ao custo de mais
desafios de step-up, ao passo que elevá-lo favorece uma experiência sem
atrito ao risco de deixar escapar ataques limítrofes. No ponto
escolhido, a matriz de confusão no conjunto de teste é dominada por
verdadeiros negativos, e os erros residuais concentram-se na fronteira
entre acessos legítimos levemente anômalos e acessos suspeitos de baixa
intensidade --- exatamente a região que o modelo de ruído foi projetado
para tornar ambígua. Isso explica tanto o F1 abaixo de 90\% quanto a
baixa taxa de falsos positivos: o ensemble é conservador, preferindo uma
detecção perdida ocasional a penalizar usuários legítimos, viés
apropriado a um front-end de autenticação, em que o atrito de
usabilidade tem custo direto de negócio. Os três modelos contribuem com
visões complementares --- a Isolation Forest sinaliza combinações de
atributos globalmente raras sem rótulos, enquanto a Random Forest e a
rede neural exploram a fronteira de decisão rotulada. Evitamos deliberadamente ler esses números como
evidência de superioridade sobre os baselines da
Tabela~\ref{tab:baseline}: os baselines não foram reexecutados sobre os
mesmos dados, e o gerador sintético, por mais cuidadosamente calibrado,
não captura toda a diversidade adversarial do tráfego de produção. Essas
observações enquadram o experimento como um estudo de viabilidade cujo
próximo passo é a validação sobre logs reais e não repudiáveis coletados
sob a própria arquitetura aqui proposta.

\textbf{Limitações e ameaças à validade.} Quatro limitações delimitam o
alcance das conclusões. Primeiro, a validade externa é restringida pelo
uso de dados sintéticos: embora o processo gerador reproduza
distribuições plausíveis de uso corporativo, ele não captura a cauda
longa de comportamentos legítimos atípicos nem a criatividade
adversarial observada em tráfego real, de modo que as métricas relatadas
devem ser lidas como um limite superior otimista e não como desempenho
esperado em produção. Segundo, a validade de construção depende da
rotulagem heurística: por derivarem das mesmas regras que descrevem os
padrões de fraude, os rótulos podem introduzir uma circularidade em que
o ensemble reaprende as fronteiras codificadas no próprio gerador ---
risco que apenas a rotulagem sobre incidentes reais, revisada por
analistas, elimina por completo. Terceiro, a comparação com os baselines
é deliberadamente qualitativa, pois estes não foram reexecutados sobre o
mesmo conjunto de dados; afirmações quantitativas de superioridade
permanecem, portanto, fora do escopo desta fase preliminar. Quarto, a
robustez adversarial é baixa: como mostra a Tabela~\ref{tab:evasion}, a
maior parte dos critérios pode ser imitada por um atacante informado.
Mitigamos parcialmente essas ameaças fixando as sementes aleatórias,
publicando os hiperparâmetros e disponibilizando abertamente o código e
os artefatos para inspeção independente; a superação plena, contudo, exige a coleta de logs
reais sob a arquitetura proposta --- registro que a camada de auditoria,
uma vez implementada, deverá tornar verificável ---, fechando o ciclo
entre a avaliação metodológica e a operação efetiva do sistema.

\section{Considerações Finais e Trabalhos Futuros}
\label{sec:conclusion}
Este artigo apresentou um sistema autônomo de autenticação contextual que
combina uma decisão adaptativa por IA com o projeto de uma auditoria
em blockchain permissionada. As contribuições são a arquitetura que
funde quatorze sinais contextuais e comportamentais por meio de um
ensemble, o projeto da camada de auditoria acoplada à decisão, com
justificativa explícita para o uso de blockchain, e uma avaliação
\emph{offline} reprodutível do motor de risco. Experimentos preliminares sobre
dados sintéticos indicaram viabilidade técnica (98,1\% de acurácia, 80,2\%
de F1-score, AUC de 0,99, latência de inferência de $\approx$33\,ms;
7,3\% de passagem silenciosa com as três faixas), e o ensemble superou
cada modelo isolado; a análise de evasão, contudo, mostra que a maior
parte dos critérios é controlável por um atacante informado. Isso
constitui uma validação metodológica e não um substituto para a
avaliação com dados reais. O próximo passo imediato é integrar uma rede
Fabric real com ao menos duas organizações, medindo latência de
\emph{commit}, vazão e detecção de adulteração, e alinhar o protótipo à
Eq.~(\ref{eq:ensemble}) e às três faixas. Como trabalhos futuros,
pretende-se ainda substituir
o conjunto sintético por dados reais ou semirreais coletados em ambiente
controlado, selecionar pesos e limiares em partição de validação,
incorporar critérios difíceis de forjar e defesas contra evasão, avaliar
desempenho e escalabilidade sob carga realista e aprofundar a análise de
segurança frente ao modelo de ameaças. Planeja-se ainda investigar a
explicabilidade das decisões do ensemble --- expondo, para cada bloqueio
ou desafio de MFA, os critérios de maior contribuição ao escore de
risco ---, avaliar o comportamento do sistema sob deriva de conceito
prolongada e ataques adversariais de evasão, e conduzir um estudo de
usabilidade que quantifique o atrito percebido pelos usuários legítimos
frente aos desafios adaptativos.

\balance
\begin{thebibliography}{00}
\bibitem{nist80063b} National Institute of Standards and Technology, ``Digital Identity Guidelines: Authentication and Lifecycle Management,'' NIST SP 800-63B, 2017.
\bibitem{verizon2024} Verizon, ``2024 Data Breach Investigations Report (DBIR),'' Verizon Business, 2024.
\bibitem{rose2020} S. Rose, O. Borchert, S. Mitchell, and S. Connelly, ``Zero Trust Architecture,'' NIST SP 800-207, 2020.
\bibitem{wiefling2019} S. Wiefling, L. Lo Iacono, and M. D\"urmuth, ``Is This Really You? An Empirical Study on Risk-Based Authentication Applied in the Wild,'' in \emph{IFIP SEC}, 2019, pp. 134--148.
\bibitem{wiefling2023} S. Wiefling, P. R. J\o rgensen, S. Thunem, and L. Lo Iacono, ``Pump Up Password Security! Evaluating and Enhancing Risk-Based Authentication on a Real-World Large-Scale Online Service,'' \emph{ACM Trans. Priv. Secur.}, vol. 26, no. 1, 2023.
\bibitem{freeman2016} D. Freeman, S. Jain, M. D\"urmuth, B. Biggio, and G. Giacinto, ``Who Are You? A Statistical Approach to Measuring User Authenticity,'' in \emph{Proc. Netw. Distrib. Syst. Secur. Symp. (NDSS)}, 2016.
\bibitem{alaca2016} F. Alaca and P. C. van Oorschot, ``Device Fingerprinting for Augmenting Web Authentication: Classification and Analysis of Methods,'' in \emph{Proc. 32nd Annu. Comput. Secur. Appl. Conf. (ACSAC)}, 2016, pp. 289--301.
\bibitem{andriamilanto2021} N. Andriamilanto, T. Allard, G. Le Guelvouit, and A. Garel, ``A Large-scale Empirical Analysis of Browser Fingerprints Properties for Web Authentication,'' \emph{ACM Trans. Web}, vol. 16, no. 1, pp. 1--62, 2021.
\bibitem{abuhamad2021} M. Abuhamad, A. Abusnaina, D. Nyang, and D. Mohaisen, ``Sensor-Based Continuous Authentication of Smartphones' Users Using Behavioral Biometrics: A Contemporary Survey,'' \emph{IEEE Internet Things J.}, vol. 8, no. 1, pp. 65--84, 2021.
\bibitem{blockaudit} A. Ahmad, M. Saad, and A. Mohaisen, ``Secure and Transparent Audit Logs with BlockAudit,'' \emph{J. Netw. Comput. Appl.}, vol. 145, art.\ 102406, 2019.
\bibitem{acien2022} A. Acien, A. Morales, R. Vera-Rodriguez, J. Fierrez, and J. V. Monaco, ``TypeNet: Deep Learning Keystroke Biometrics,'' \emph{IEEE Trans. Biom. Behav. Identity Sci.}, vol. 4, no. 1, pp. 57--70, 2022.
\bibitem{stragapede2024} G. Stragapede, P. Delgado-Santos, R. Tolosana, R. Vera-Rodriguez, R. Guest, and A. Morales, ``TypeFormer: Transformers for Mobile Keystroke Biometrics,'' \emph{Neural Comput. Appl.}, vol. 36, pp. 18531--18545, 2024.
\bibitem{frank2013} M. Frank, R. Biedert, E. Ma, I. Martinovic, and D. Song, ``Touchalytics: On the Applicability of Touchscreen Input as a Behavioral Biometric for Continuous Authentication,'' \emph{IEEE Trans. Inf. Forensics Security}, vol. 8, no. 1, pp. 136--148, 2013.
\bibitem{liu2008} F. T. Liu, K. M. Ting, and Z.-H. Zhou, ``Isolation Forest,'' in \emph{IEEE Int. Conf. Data Mining (ICDM)}, 2008, pp. 413--422.
\bibitem{killourhy2009} K. S. Killourhy and R. A. Maxion, ``Comparing Anomaly-Detection Algorithms for Keystroke Dynamics,'' in \emph{IEEE/IFIP DSN}, 2009, pp. 125--134.
\bibitem{goodfellow2016} I. Goodfellow, Y. Bengio, and A. Courville, \emph{Deep Learning}. Cambridge, MA: MIT Press, 2016.
\bibitem{gama2014} J. Gama, I. \v{Z}liobait\.e, A. Bifet, M. Pechenizkiy, and A. Bouchachia, ``A Survey on Concept Drift Adaptation,'' \emph{ACM Comput. Surv.}, vol. 46, no. 4, pp. 1--37, 2014.
\bibitem{fabric2024} Hyperledger Foundation, ``Hyperledger Fabric Documentation, v2.5,'' Linux Foundation, 2024.
\bibitem{nistir8202} National Institute of Standards and Technology, ``Blockchain Technology Overview,'' NISTIR 8202, 2018.
\bibitem{nist80053r5} National Institute of Standards and Technology, ``Security and Privacy Controls for Information Systems and Organizations,'' NIST SP 800-53 Rev.\ 5, 2020.
\bibitem{lgpd2018} Brasil, ``Lei n.\ 13.709/2018 -- Lei Geral de Proteção de Dados Pessoais (LGPD),'' 2018.
\end{thebibliography}

\end{document}
